\subsection*{Signatures}
\label{sec:axiom-systems-signatures}

\begin{toolkitbox}{Signatures Toolkit --- Quick Reference}
\small
\begin{tabular}{l l l}
\toprule
\textbf{Concept} & \textbf{Role} & \textbf{Detail} \\
\midrule
Signature & non-logical vocabulary & \hyperref[def:signature]{Definition} \\
Language & signature together with logical syntax & \hyperref[def:language]{Definition} \\
Constant symbol & names a distinguished object syntactically & \hyperref[def:signature-constant-symbol]{Definition} \\
Function symbol & names an operation syntactically & \hyperref[def:signature-function-symbol]{Definition} \\
Relation symbol & names a relation syntactically & \hyperref[def:relation-symbol]{Definition} \\
Arity & records number of input places & \hyperref[def:arity]{Definition} \\
\bottomrule
\end{tabular}
\end{toolkitbox}

\begin{exposition}
\SignatureMetadata{arithmetic-signature}
\end{exposition}

\begin{signaturebox}{Signature Presentation: A Simple Arithmetic Language}
Let
\[
  \Sigma_{\mathrm{arith}}
  =
  \bigl(
    \mathcal{C}_{\mathrm{arith}},
    \mathcal{F}_{\mathrm{arith}},
    \mathcal{R}_{\mathrm{arith}},
    \operatorname{ar}
  \bigr)
\]
where
\[
  \mathcal{C}_{\mathrm{arith}}=\{0,1\},\qquad
  \mathcal{F}_{\mathrm{arith}}=\{+,\cdot,S\},\qquad
  \mathcal{R}_{\mathrm{arith}}=\{<\}.
\]
The arity function is specified by
\[
  \operatorname{ar}(0)=\operatorname{ar}(1)=0,\qquad
  \operatorname{ar}(S)=1,\qquad
  \operatorname{ar}(+)=\operatorname{ar}(\cdot)=2,\qquad
  \operatorname{ar}(<)=2.
\]

\begin{signaturetable}
$0,1$ & constants & $0$ & distinguished constant names \\
$S$ & function symbol & $1$ & successor operation symbol \\
$+,\cdot$ & function symbols & $2$ & binary operation symbols \\
$<$ & relation symbol & $2$ & binary order relation symbol \\
\end{signaturetable}

The signature records only the available non-logical symbols and their arities.
It does not assert any axioms about those symbols.
\end{signaturebox}

\begin{definition}[Signature]
\label{def:signature}
A \emph{signature} is a collection of non-logical symbols divided into
constant symbols, function symbols, and relation symbols, together with an
arity assignment for each symbol. Constant symbols are recorded as arity-$0$
symbols.
\end{definition}

\begin{remark*}[Interpretation]
A signature is the formal vocabulary available before any axioms are chosen.
It says which symbols may occur and how many input places each symbol has.
\end{remark*}

\NoLocalDependencies

\begin{definition}[Arity]
\label{def:arity}
The \emph{arity} of a symbol is the number of input places assigned to that
symbol by a signature. Constants have arity $0$; unary symbols have arity $1$;
binary symbols have arity $2$.
\end{definition}

\begin{remark*}[Interpretation]
Arity is bookkeeping for grammatical well-formedness. It determines whether a
symbol has been supplied with the correct number of arguments.
\end{remark*}

\begin{dependencies}
\begin{itemize}
  \item \hyperref[def:signature]{Signature}
\end{itemize}
\end{dependencies}

\begin{definition}[Language]
\label{def:language}
A \emph{language} determined by a signature $\Sigma$ consists of the
non-logical symbols of $\Sigma$ together with the logical symbols, variables,
punctuation, and formation rules used to build terms and formulas.
\end{definition}

\begin{remark*}[Interpretation]
The signature supplies the subject-specific vocabulary. The language adds the
logical grammar needed to form expressions from that vocabulary.
\end{remark*}

\begin{dependencies}
\begin{itemize}
  \item \hyperref[def:signature]{Signature}
\end{itemize}
\end{dependencies}

\begin{definition}[Constant Symbol]
\label{def:signature-constant-symbol}
A \emph{constant symbol} is a non-logical symbol of arity $0$ in a signature.
It may occur as a term without taking any input terms.
\end{definition}

\begin{remark*}[Interpretation]
A constant symbol is a name-forming symbol. In the arithmetic signature above,
$0$ and $1$ are constant symbols.
\end{remark*}

\begin{dependencies}
\begin{itemize}
  \item \hyperref[def:signature]{Signature}
  \item \hyperref[def:arity]{Arity}
\end{itemize}
\end{dependencies}

\begin{definition}[Function Symbol]
\label{def:signature-function-symbol}
A \emph{function symbol} is a non-logical symbol in a signature assigned an
arity $n\geq 1$. If $f$ has arity $n$ and $t_1,\ldots,t_n$ are terms, then
$f(t_1,\ldots,t_n)$ is a term.
\end{definition}

\begin{remark*}[Interpretation]
A function symbol is an operation-forming symbol. In the arithmetic signature
above, $S$ is unary, while $+$ and $\cdot$ are binary.
\end{remark*}

\begin{dependencies}
\begin{itemize}
  \item \hyperref[def:signature]{Signature}
  \item \hyperref[def:arity]{Arity}
\end{itemize}
\end{dependencies}

\begin{definition}[Relation Symbol]
\label{def:relation-symbol}
A \emph{relation symbol} is a non-logical symbol in a signature assigned an
arity $n\geq 1$. If $R$ has arity $n$ and $t_1,\ldots,t_n$ are terms, then
$R(t_1,\ldots,t_n)$ is an atomic formula.
\end{definition}

\begin{remark*}[Interpretation]
A relation symbol is a statement-forming symbol once it is applied to the
right number of terms. In the arithmetic signature above, $<$ is a binary
relation symbol.
\end{remark*}

\begin{dependencies}
\begin{itemize}
  \item \hyperref[def:signature]{Signature}
  \item \hyperref[def:arity]{Arity}
\end{itemize}
\end{dependencies}

\begin{remark*}[Structural remarks]
The same signature may support different choices of axioms later in the
chapter. For example, a single arithmetic vocabulary can be used with stronger
or weaker axiom sets.

The same signature may also be read over different mathematical universes once
interpretations are introduced. At this stage, the signature itself records
only the common vocabulary; it does not choose any particular interpretation.
\end{remark*}
